\documentclass[10pt,a4paper]{amsart}
\usepackage[margin=19mm]{geometry}
\usepackage{amsmath,amssymb,mathtools}
\usepackage[scr=rsfs,cal=euler]{mathalfa}
\usepackage{xfrac}
\usepackage[unicode,hidelinks,bookmarks=false]{hyperref}
\newcommand{\ie}{i.e.\ }
\newcommand{\cf}{cf.\ }
\newcommand{\resp}{respectively\ }
\newcommand{\Subsection}{\subsection}
\providecommand{\nobreakdash}{\nobreak\hbox{-}}
\setlength{\emergencystretch}{2em}
\multlinegap0pt
\usepackage{bm}
\usepackage{bbm}
\usepackage{dsfont}
\usepackage{xspace}
\usepackage{cancel}
\usepackage{mathtools}
\usepackage{amsthm}
\makeatletter
\@ifundefined{theorem}{\newtheorem{theorem}{Theorem}}{}
\makeatother
\title[Explicit Density Approximation for Neural Implicit Samplers ]{Explicit Density Approximation for Neural Implicit Samplers Using a Bernstein-Based Convex Divergence}
\author{Jos\'e Manuel de Frutos, Manuel A. V\'azquez, Pablo M. Olmos, Joaqu\'in M\'iguez}
\date{Source: arXiv:2506.04700v2 (CC BY 4.0)}
\begin{document}
\maketitle

\noindent\emph{Source and licence.} Explicit Density Approximation for Neural Implicit Samplers Using a Bernstein-Based Convex Divergence. Version arXiv:2506.04700v2 (CC BY 4.0), distributed under the Creative Commons Attribution 4.0 International licence, as recorded in the arXiv licence metadata for this exact version and shown on the abstract page https://arxiv.org/abs/2506.04700v2. Full rights notice: licenses/arxiv-2506.04700-CC-BY-4.0.txt.\par\smallskip

\noindent\textit{Editorial scope.} Counted unit: the Theorem whose source label is \texttt{theorem: Finite-Sample Deviation Bound} in the article (source0.tex lines 1761--1823, source label \texttt{theorem: Finite-Sample Deviation Bound}), delivered together with the article's own complete original proof of it (source0.tex lines 1825--2010, one \texttt{proof} environment of 186 lines). No mathematics is written, repaired or re-derived: statement and proof are the article's own text line for line. Required context retained uncounted: none. Every definition, display and label of those lines is retained unchanged. Omitted from this package and not redistributed: 1--1760, 1824--1824, 2011--4180. Nothing inside a retained excerpt is omitted or altered: every retained body is delivered exactly as the article's lines are written, so no omission receipt of any kind accompanies this package. Citations: the retained excerpts carry no citation command at all, so no bibliography entry of the article is transcribed and no reference is redistributed. Reference adaptation: none was needed and none was made; every reference inside the delivered unit resolves inside it, so no unresolved reference or citation is produced. Printed slips kept literal: no printed word, symbol, hypothesis or number of the counted statement or of its proof was altered. Layout and class: the article's own class is \texttt{article} with packages including neurips\_2024, biblatex, inputenc, fontenc, hyperref, url, booktabs, pgfplots, multirow, amsfonts, nicefrac, microtype, amsmath, amsthm; the wrapper is the lane's pinned \texttt{amsart} editorial wrapper at 10pt on A4, which restates the theorem-like environments the retained text needs and adds the article's own macro definitions that the retained text needs (none), with extra wrapper packages (bm, bbm, dsfont, xspace, cancel, mathtools). The delivered numbering is the unit's own. Assets: no retained span contains a figure environment, a graphics-inclusion command, a PDF-page inclusion, a table or a TikZ picture; no image file is copied into this package, the package declares no asset dependency and no image file is redistributed with it. Licence: version arXiv:2506.04700v2 of this article is distributed under the Creative Commons Attribution 4.0 International licence, as recorded in the arXiv licence metadata for this exact version and shown on the abstract page https://arxiv.org/abs/2506.04700v2. A full rights notice is delivered at licenses/arxiv-2506.04700-CC-BY-4.0.txt. Characters: every retained excerpt is pure ASCII, so the delivered engine, XeLaTeX with its default Latin Modern text font, reports no missing character.
\begin{theorem}[Finite-Sample Deviation Bound for the Rank-Based Discrepancy] \label{theorem: Finite-Sample Deviation Bound}
\label{thm:rank_dev_bound}
Let $p$ be a continuous pdf on a compact set $\mathcal{K} \subseteq \mathbb{R}$, 
and let $\{X_i\}_{i=1}^n$ be $n$ i.i.d.\ samples from $p$. 
Define the empirical measure $\hat{p}_n$ to assign mass $\tfrac{1}{n}$ to each sample point $X_i$. 
Let $\tilde{p}$ be another pdf on $\mathcal{K}$ with cdf $\tilde{F}$. 
Consider the rank-based discrepancy
\[
d_{K}(p,\tilde{p}) 
~:=~
\frac{1}{K+1} 
\sum_{r=0}^{K}
\Bigl|\,
  \tfrac{1}{K+1}
  ~-~
  \mathbb{Q}_{K}(r) 
\Bigr|,
\]
where
\[
\mathbb{Q}_{K}(r)
~:=~
\int_{\mathcal{K}}
  \binom{K}{r}\,
  \bigl[\tilde{F}(x)\bigr]^r
  \,\bigl[1 - \tilde{F}(x)\bigr]^{K-r}
\,p(\mathrm{d}x).
\]
Likewise, define its empirical analog using $\hat{p}_n$:
\[
\hat{d}_{K}\bigl(\hat{p}_n,\tilde{p}\bigr)
~:=~
\frac{1}{K+1}
\sum_{r=0}^{K}
\Bigl|\,
  \tfrac{1}{K+1}
  ~-~
  \hat{\mathbb{Q}}_{K}(r)
\Bigr|,
\quad
\text{where}
\quad
\hat{\mathbb{Q}}_{K}(r)
~:=~
\int_{\mathcal{K}}
  \binom{K}{r}\,
  [\tilde{F}(x)]^r
  \,[\,1 - \tilde{F}(x)\bigr]^{K-r}
\,\hat{p}_n(\mathrm{d}x).
\]

Then, for any $\delta > 0$, with probability at least $1 - \delta$ (over the i.i.d.\ draw of $\{X_i\}$), we have
\[
\bigl|\,
  \hat{d}_{K}(\hat{p}_n, \tilde{p})
  ~-~
  d_{K}(p, \tilde{p})
\bigr|
~\le~
C \sqrt{\frac{\log\!\bigl(\tfrac{1}{\delta}\bigr)}{n}},
\]
where $C$ is a constant that depends on $K$ (and can often be taken on the order of $(K+1)$, or possibly smaller by tighter analysis).
\end{theorem}

\begin{proof}[Proof of Theorem \ref{theorem: Finite-Sample Deviation Bound}]
Recall
\begin{align*}
\hat{\mathbb{Q}}_K(r)
=
\int_{\mathbb{R}}
  \binom{K}{r}\,
  \tilde{F}(x)^r(1 - \tilde{F}(x)\bigr)^{K-r}
\,\hat{p}_n(x)\mathrm{d}x
=
\frac{1}{n} \sum_{i=1}^{n}
  \binom{K}{r}
  \tilde{F}(X_i)^r
  (1-\tilde{F}(X_i))^{K-r},
\end{align*}
while
\begin{align*}
\mathbb{Q}_{K}(r)=
\int_{\mathbb{R}}
  \binom{K}{r}\,
  \tilde{F}(x)^{r}(1-\tilde{F}(x))^{K-r}
\,p(x)\mathrm{d}x
=
\mathbb{E}_{X \sim p}\Bigl[
  \binom{K}{r}\,\tilde{F}(X)^r(1-\tilde{F}(X))^{K-r}
\Bigr].
\end{align*}

For each $r \in \{0,\ldots,K\}$, note that 
\begin{align*}
\binom{K}{r}\tilde{F}(x)^r (1 - \tilde{F}(x))^{K-r}
~\le~
1,
\end{align*}
since $\tilde{F}(x)\in [0,1]$ and $\binom{K}{r}$ is precisely the binomial coefficient in $[\tilde{F}(x)]^r[1-\tilde{F}(x)]^{K-r}$.  
Hence each random variable
\begin{align*}
Z_i^{(r)}
~:=~
\binom{K}{r}\,
[\tilde{F}(X_i)]^r\,
[1 - \tilde{F}(X_i)]^{K-r}
\end{align*}
is in $[0,1]$ a.s.

By Hoeffding's inequality (or a simple Chernoff bound for bounded variables):
\begin{align*}
\mathbb{P}\!\Bigl[
  \bigl|\hat{\mathbb{Q}}_K(r) - \mathbb{Q}_K(r)\bigr|
  ~>~
  \epsilon
\Bigr]
~\le~
2\exp\Bigl(
  -2\,n\,\epsilon^2
\Bigr).
\end{align*}

Define
\begin{align*}
\Delta(r)=
\bigl|\,
  \hat{\mathbb{Q}}_K(r)
  ~-~
  \mathbb{Q}_K(r)
\bigr|.
\end{align*}
We want all $\Delta(r)\le \epsilon$. By the union bound,
\begin{align*}
\mathbb{P}\!\Bigl[
  \exists\,r\in \{0,\ldots,K\}:\,\Delta(r) > \epsilon
\Bigr]
~\le~
(K+1)\,\cdot 2\exp(-2n\epsilon^2).
\end{align*}
Hence with probability at least $1-(K+1)2\exp(-2n\epsilon^2)$, we have $\Delta(r)\le \epsilon$ for \emph{all} $r$.  

Recall
\begin{align*}
d_{K}(p,\tilde{p})
=\frac{1}{K+1}
  \sum_{r=0}^K
  \Bigl|
    \tfrac{1}{K+1} - \mathbb{Q}_{K}(r)
  \Bigr|,
\end{align*}
and
\begin{align*}
\hat{d}_{K}(\hat{p}_n,\tilde{p})
=\frac{1}{K+1}
  \sum_{r=0}^K
  \Bigl|
    \tfrac{1}{K+1} - \hat{\mathbb{Q}}_{K}(r)
  \Bigr|.
\end{align*}
Hence by triangle inequality
\begin{align*}
\bigl|
  \hat{d}_{K}(\hat{p}_n,\tilde{p}) 
  -
  d_{K}(p,\tilde{p})
\bigr|\le
\frac{1}{K+1}
\sum_{r=0}^{K}
\left|
  \bigl|
    \tfrac{1}{K+1} - \hat{\mathbb{Q}}_{K}(r)
  \bigr|
  -
  \bigl|
    \tfrac{1}{K+1} - \mathbb{Q}_{K}(r)
  \bigr|
\right|.
\end{align*}
Using the standard inequality
$\bigl||a|-|b|\bigr|\le |a-b|,$ 
one obtains
\begin{align*}
\bigl|
  \hat{d}_{K}(\hat{p}_n,\tilde{p})
  -
  d_{K}(p,\tilde{p})
\bigr|
&\le
\frac{1}{K+1}
\sum_{r=0}^K
\Bigl|
  \hat{\mathbb{Q}}_{K}(r)
  -
  \mathbb{Q}_{K}(r)
\Bigr|
\\&\le
\frac{K+1}{K+1}\,\max_{r}\Delta(r)
=
\max_{r}\Delta(r).
\end{align*}
Thus
\begin{align*}
\bigl|
  \hat{d}_{K}(\hat{p}_n,\tilde{p}) 
  -
  d_{K}(p,\tilde{p})
\bigr|
\le
\epsilon
\quad
\text{whenever }
\Delta(r)\le \epsilon ~\forall\,r.
\end{align*}

Putting it all together,
\begin{align*}
\mathbb{P}\Bigl[
  \bigl|\hat{d}_{K}(\hat{p}_n,\tilde{p}) 
  -
  d_{K}(p,\tilde{p})\bigr|
  ~>~
  \epsilon
\Bigr]
~\le~
\mathbb{P}\Bigl[
  \exists\,r\text{ s.t. } \Delta(r)>\epsilon
\Bigr]
~\le~
(K+1)\cdot 2\,\exp(-2n\epsilon^2).
\end{align*}
Solving for $\epsilon$ in terms of $\delta=(K+1)2\,\exp(-2n\epsilon^2)$ yields 
\begin{align*}
\epsilon
~\lesssim~
\sqrt{\frac{\log\!\bigl(\tfrac{(K+1)2}{\delta}\bigr)}{2n}}
,
\end{align*}
and we absorb constants to get a final form 
\begin{align*}
\bigl|
  \hat{d}_{K}(\hat{p}_n,\tilde{p})
  -
  d_{K}(p,\tilde{p})
\bigr|
\le
C
\,\sqrt{\frac{\log\!\bigl(\tfrac{1}{\delta}\bigr)}{n}}.
\end{align*}
The constant $C$ can be chosen to handle the factor $(K+1)$ inside the log, etc., typically $C$ might be 2 or $2\sqrt{2}$ depending on details.
\end{proof}

\end{document}
